% C3 foundational module. This is an ORIGINAL invariant-group-action result.
% It has no dependency on terminal Ramsey, typed strong profiles or SAP.
\section{Invariant typed quotient transport and canonical return holonomy}
\label{sec:foundational-typed-return-holonomy}

\paragraph{Invariant typed quotient transport (foundational).}
The next lemma needs only original named transitivity and invariance of
the represented binary relation; neither SAP nor CCO nor exactness
of that relation is required. We write $1_a$ for the FULL universal
same-sort relation, as distinct from logical identity $\Id_a$.
Transport preserves the interval of ALL invariant equivalences.
Whether a transported invariant equivalence is also FULL RAW exact
is a separate assertion and is not inferred from interval transport.

\begin{lemma}[Typed quotient transport]\label{lem:typed-quotient-transport}
Let \(R\) be a nonempty invariant binary relation from a named original transitive sort \(a\) to a named original transitive sort \(b\), and let \(E\) be an invariant equivalence on sort \(a\). Suppose
\[
 RR^{\smile}\subseteq E.
\tag{QT1}\label{eq:typed-transport-kernel}
\]
For \(y\in b\), let \(\lambda_{R,E}(y)\) be the \(E\)-class of any
\(x\) with \(xRy\), and put
\[
 E^R=R^{\smile}ER.
\]
Then \(\lambda_{R,E}\) is well defined, \(E^R\) is a definable equivalence,
and \(R\) induces a bijection
\[
 q_{R,E}:a/E\longrightarrow b/E^R.
\]
Transport by this bijection is a lattice isomorphism
\[
 [E,1_a]\longrightarrow [E^R,1_b],
 \qquad H\longmapsto H^R=R^{\smile}HR.
\tag{QT2}\label{eq:typed-transport-lattice}
\]
It respects composition of composable relations for actual composable nonempty invariant typed relations whenever the displayed kernel hypotheses hold. In particular,
meet-irreducibility is preserved by transport.
An ADDITIONAL noncollapse conclusion under visible SAP
may hold in appropriate terminal sources, but is NOT part
of this general invariant quotient-transport assertion.
\end{lemma}

\begin{proof}
Since the original endpoint sorts are transitive,
the nonempty invariant relation $R$ is total and
surjective. For $y\in b$ its nonempty set of
$R$-predecessors lies in ONE $E$-class, because
$RR^\smile\subseteq E$. This gives a surjective
equivariant map $\lambda_{R,E}:b\to a/E$.
For any $y,y'\in b$,
\[
 \lambda_{R,E}(y)=\lambda_{R,E}(y')
  \quad\Longleftrightarrow\quad
             y(R^\smile;E;R)y',
\]
since any two predecessors of either point
are already $E$-equivalent. Hence $E^R$
is precisely the kernel of $\lambda_{R,E}$;
the quotient map induced by $\lambda_{R,E}$
is a bijection $b/E^R\to a/E$, whose inverse
is the displayed $q_{R,E}:a/E\to b/E^R$.

For any invariant $H\ge E$ on $a$, the
condition $y(R^\smile;H;R)y'$ holds exactly
when $\lambda_{R,E}(y)$ and
$\lambda_{R,E}(y')$ are $H/E$-equivalent.
Thus conjugation gives the full interval
isomorphism and preserves covers and
meet-irreducibles. If $R:a\to b$, $S:b\to c$ satisfy
the corresponding two quotient kernel
conditions and $F=E^R$, uniqueness of
the E-class of each $(R;S)$-predecessor
gives directly
\[
 q_{R;S,E}=q_{S,F}\circ q_{R,E},
 \qquad E^{R;S}=F^S.
\]
Thus composable transports respect their
actual RAW relational composition.
No mixed-sort commutativity is assumed.
\end{proof}

\begin{lemma}[Canonical invariant typed return and identity holonomy]\label{lem:foundational-typed-return-holonomy}
Let \(P_1,\ldots,P_n\) be a nonempty finite composable sequence
of nonempty invariant represented relations between named transitive
original sorts, from \(a_0\) back to \(a_0\), and suppose
\[
 P_1\cdots P_n\subseteq E_0
\]
for a definable equivalence \(E_0\) on \(a_0\).  Put
\(R_i=P_1\cdots P_i\).  Then every prefix transports \(E_0\) to a definable
equivalence \(E_i\) on the intermediate sort, every factor is
quotient-functional between consecutive transported quotients, and the whole
loop has identity holonomy on \(a_0/E_0\).  The same conclusion holds
simultaneously for EVERY invariant equivalence \(J\supseteq E_0\).
No representation of the transported \(J_i\) by named FULL RAW exact
quotient sorts is asserted.
\end{lemma}

\begin{proof}
We make the quotient labels at EVERY intermediate position
CANONICAL; this avoids choosing unrelated transport charts.
Write \(a_i\) for the original real sort after \(P_i\),
with \(a_n=a_0\), and put
\[
 R_0=\Id_{a_0},\quad
 R_i=P_1;\cdots;P_i,\qquad
 S_i=P_{i+1};\cdots;P_n,\quad S_n=\Id_{a_0}.
\]
All products here are ACTUAL equality-inclusive RAW represented
relations. Since the \(P_i\) are nonempty invariant relations between
TRANSITIVE named original sorts, they and all composable prefixes
and suffixes are serial and co-serial.

For each \(i\) and real \(y\in a_i\), choose ANY predecessor
\(x\in a_0\) with \(xR_i y\), and define
\[
             \lambda_i(y)=[x]_{E_0}\in a_0/E_0.
\]
The value is independent of \(x\): if also \(x'R_i y\),
choose a suffix witness \(z\) with \(yS_i z\).
Then \(xR_nz\) and \(x'R_nz\), and the hypothesis
\(R_n\subseteq E_0\) gives \(xE_0zE_0x'\).
Co-seriality supplies a predecessor for every \(y\);
seriality of \(R_i\) shows that every \(E_0\)-class
occurs as some value of \(\lambda_i\).
In particular this is ONE onto original quotient labelling,
not an independent choice at each edge.

Put \(E_i=\ker(\lambda_i)\). The witness characterization
gives the literal equality of invariant ORIGINAL equivalences
\[
                 E_i=R_i^{\smile};E_0;R_i .
\]
Thus the preceding typed quotient-transport lemma applies to
each prefix \(R_i\), including all repeated original sorts.
For any actual \(yP_i z\), choose \(xR_{i-1}y\).
Then \(xR_i z\), so
\[
                   \lambda_{i-1}(y)=\lambda_i(z).
\]
It follows directly that \(P_iP_i^{\smile}\subseteq E_{i-1}\)
and \(P_i^{\smile};E_{i-1};P_i=E_i\).
The induced quotient map
\(q_i:a_{i-1}/E_{i-1}\to a_i/E_i\)
is therefore a BIJECTION.  More precisely, the
bijections \(\bar\lambda_i:a_i/E_i\to a_0/E_0\)
induced by \(\lambda_i\) satisfy
\[
           q_i=\bar\lambda_i^{-1}\circ\bar\lambda_{i-1}.
\]
At the start \(\lambda_0(x)=[x]_{E_0}\), since \(R_0=\Id\).
At the end every predecessor \(xR_ny\) satisfies
\(xE_0y\) by the whole-word hypothesis, so
\(\lambda_n(y)=[y]_{E_0}\) on the SAME endpoint sort.
Hence
\[
            q_n\circ\cdots\circ q_1
               =\bar\lambda_n^{-1}\circ\bar\lambda_0
               =\operatorname{id}_{a_0/E_0}.
\]
This proves IDENTITY holonomy without selecting a
groupoid section, an origin system, or a quotient chart.

Finally, for EVERY equivalence \(J\supseteq E_0\),
compose all the preceding labels with the canonical
projection \(a_0/E_0\to a_0/J\).
The new canonical kernels are SIMULTANEOUSLY
\[
              J_i=R_i^{\smile};J;R_i,
\]
and the same edge and end-of-return identities hold on
\(a_0/J\).
All these kernels and maps are ACTUAL invariant quotient objects.
This abstract conclusion uses neither SAP nor the represented exact
lattice, and does not claim that any intermediate \(J_i\) is a
named unary-coordinate sort of a later expansion.
\end{proof}

\paragraph{Prescribed quotient actions on points.}
The following elementary action criterion will be used for the finite
owner systems, independently of any exactness or completion hypothesis.

\begin{lemma}[Prescribed-action point lifting]
\label{lem:prescribed-action-point-lifting}
Let $K$ act on a nonempty set $Y$, let $\Gamma$ act on $B$, and
let $\rho:K\to\Gamma$ be a homomorphism. Suppose $\pi:Y\to B$
is a surjection with $\pi(ky)=\rho(k)\pi(y)$. The following are
equivalent:
\begin{enumerate}[label=(\roman*)]
\item $\rho$ is onto and $\ker\rho$ is transitive on each fibre
      of $\pi$;
\item for every $x,y\in Y$ and every $\gamma\in\Gamma$ with
      $\gamma\pi(x)=\pi(y)$, some $k\in K$ satisfies
      $\rho(k)=\gamma$ and $kx=y$.
\end{enumerate}
Neither action is required to be faithful, and no group section is
assumed.
\end{lemma}
\begin{proof}
Assume (i). Lift $\gamma$ to $k_0\in K$. The points $k_0x,y$
lie in the same $\pi$-fibre, so some $h\in\ker\rho$ sends
$k_0x$ to $y$. Then $k=hk_0$ satisfies (ii).
Conversely, for any $\gamma\in\Gamma$ choose $x\in Y$ and,
by surjectivity, $y$ above $\gamma\pi(x)$. Condition (ii)
provides a lift of $\gamma$, so $\rho$ is onto. Taking
$\gamma=1$ in (ii) gives transitivity of $\ker\rho$ on each
fibre. Nonemptiness is used when choosing $x$.
\end{proof}

\paragraph{Orbital pullbacks and the real diagonal.}
The following group-action statement separates representative saturation
from any claim about positive words or STRONG reachability.

\begin{lemma}[Orbital pullbacks through equivariant quotient maps]
\label{lem:foundational-orbital-pullback}
Let $G$ act transitively on nonempty named root sorts $A,B$, which
are either the same sort or disjoint. Let $\tau:A\to Q$ and
$\sigma:B\to R$ be equivariant surjections to $G$-sets, with kernels
$J=\ker\tau$ and $K=\ker\sigma$. For a nonempty $G$-orbital
$\Theta\subseteq Q\times R$, put
\[
 P_\Theta=\{(x,y)\in A\times B:(\tau(x),\sigma(y))\in\Theta\},
 \qquad D_{A,B}=\{(x,y)\in A\times B:x\ne y\}.
\]
For any $(x_0,y_0)\in P_\Theta$, let $O=G\cdot(x_0,y_0)$.
Then, with actual equality-inclusive relational composition,
\[
                  P_\Theta=J;O;K.
\]
Thus whenever a distinct-root witness is available, the diversity
pullback is exactly $(J;O;K)\cap D_{A,B}$ for its root orbital $O$.
The nonempty diversity pullbacks partition $D_{A,B}$ and induce a
partition of the set of its root $G$-orbitals.

The diversity pullback is empty if and only if $A=B$,
$J=K=\Id_A$, and
\[
                \Theta=\{(\tau(x),\sigma(x)):x\in A\}.
\]
Hence at most one output orbital has empty diversity pullback. In
this case both maps are equivariant bijections and the missing
diversity comes from the same original root, even if $Q,R$ are
different output sorts. No homogeneity, SAP, CCO or word-realisation
assumption is needed.
\end{lemma}
\begin{proof}
Surjectivity makes $P_\Theta$ nonempty. If $(x,y)\in P_\Theta$,
some $g\in G$ sends $(\tau(x_0),\sigma(y_0))$ to
$(\tau(x),\sigma(y))$. Equivariance gives
$x\,J\,gx_0$ and $gy_0\,K\,y$, so $(x,y)\in J;O;K$.
Conversely changing the two representatives of a pair in $O$ within
its $J$- and $K$-classes leaves its output values unchanged. Those
values lie in $\Theta$, proving the equality. Equivariance assigns
each root orbital to exactly one output orbital. The asserted
partitions follow by restricting to the invariant set $D_{A,B}$.

If $P_\Theta$ has no distinct-root pair, its nonemptiness forces
$A=B$, and it contains some $(x,x)$. If $J\ne\Id_A$, transitivity
makes every $J$-class have at least two elements. Choose $x'\ne x$
with $x'Jx$; then $(x',x)\in P_\Theta$, a contradiction. The same
argument applies to $K$, so both kernels are identity. Transitivity
now makes the displayed same-point graph exactly the output orbital
of $(\tau(x),\sigma(x))$, hence exactly $\Theta$. Conversely, for
injective $\tau,\sigma$ the preimage of this graph consists precisely
of the diagonal of $A$. The graph is uniquely determined by the two
maps, proving uniqueness of the empty-diversity exception.
\end{proof}
